\section{Proof Methods [1.7-1.9]}

\subsection{Proof by Cases}

You can break a complicated proof into cases, and prove each separately. You can also break cases into subcases.

\subsection{Proof by Contradiction}

In proof by contradiction, or indirect proof, you show that if a proposition were false, then a false fact would be true. Since a false fact cannot be true, the proposition must be true.

\vspace{1mm}

To prove a proposition \(P\) by contradiction:

1) Write 'We use proof by contradiction'

2) Write 'Suppose \(P\) is false'

3) Deduce something known to be false (a logical contradiction)

4) Write 'This is a contradiction. Therefore \(P\) must be true.'

\vspace{1mm}

To prove a surd is irrational, we can use contradiction, by saying rational numbers can be expressed as \(\frac{p}{q}\) where \(p\) and \(q\) are in there lowest terms.

We can square both sides, and deduce that the fraction is not in its lowest terms, contradicting the assumption, proving the surd is irrational.

\section{Good Proofs in Practice}

A good proof must be clear. Correctness and clarity work together, as a well-written proof is more likely to be a correct proof, as mistakes are harder to hide.

Tips:
\begin{itemize}
    \item State what you are going to do
    \item Keep a linear, intelligible order
    \item Explain your calculations
    \item Avoid excessive symbolism
    \item Simplify your proof
    \item Introduce notation thoughtfully, and define their meaning
    \item Structure long proofs
    \item Be wary of the obvious facts
    \item Finish, draw everything together and show why the original claim follows
\end{itemize}

